\documentclass[12pt,a4paper]{article}

% ── Codificação e idioma ─────────────────────────────────────────────
\usepackage[utf8]{inputenc}
\usepackage[T1]{fontenc}
\usepackage[portuguese]{babel}
\usepackage{lmodern}
\usepackage{microtype}

% ── Layout ───────────────────────────────────────────────────────────
\usepackage[top=2.8cm,bottom=2.8cm,left=3.0cm,right=2.5cm,headheight=16pt]{geometry}

% ── Cores ────────────────────────────────────────────────────────────
\usepackage[dvipsnames,table]{xcolor}
\definecolor{navy}{HTML}{1A3A5C}
\definecolor{medblue}{HTML}{2E75B6}
\definecolor{ltblue}{HTML}{D6E4F0}
\definecolor{rowalt}{HTML}{EBF3FB}
\definecolor{legalbg}{HTML}{F0FDF4}
\definecolor{legalbdr}{HTML}{16A34A}
\definecolor{alertbg}{HTML}{FFF8E1}
\definecolor{alertbdr}{HTML}{E59C0A}
\definecolor{warnbg}{HTML}{FEF2F2}
\definecolor{warnbdr}{HTML}{DC2626}
\definecolor{infobg}{HTML}{EFF6FF}
\definecolor{infobdr}{HTML}{3B82F6}
\definecolor{grayrow}{HTML}{F3F4F6}

% ── Cabeçalho / rodapé ───────────────────────────────────────────────
\usepackage{fancyhdr,lastpage}
\pagestyle{fancy}\fancyhf{}
\fancyhead[L]{\footnotesize\color{navy}\textbf{TR} --- Materiais de Limpeza e Higiene --- Borba/AM}
\fancyhead[R]{\footnotesize\color{navy}Lei n.\textsuperscript{o} 14.133/2021}
\fancyfoot[C]{\footnotesize Página \thepage\ de \pageref{LastPage}}
\renewcommand{\headrulewidth}{0.5pt}
\renewcommand{\headrule}{\hbox to\headwidth{\color{medblue}\leaders\hrule height 0.5pt\hfill}}

% ── Caixas coloridas ─────────────────────────────────────────────────
\usepackage{tcolorbox}
\tcbuselibrary{skins,breakable}

\newtcolorbox{legalbox}[1][Base Legal]{enhanced,breakable,
  colback=legalbg,colframe=legalbdr,boxrule=0.8pt,arc=4pt,
  left=10pt,right=10pt,top=5pt,bottom=5pt,
  title=\textbf{#1},fonttitle=\small\bfseries\color{legalbdr},
  attach boxed title to top left={yshift=-2mm,xshift=6mm},
  boxed title style={colback=legalbg,colframe=legalbdr,arc=3pt,boxrule=0.6pt}}

\newtcolorbox{alertbox}[1][Atenção]{enhanced,breakable,
  colback=alertbg,colframe=alertbdr,boxrule=0.8pt,arc=4pt,
  left=10pt,right=10pt,top=5pt,bottom=5pt,
  title=\textbf{#1},fonttitle=\small\bfseries\color{navy},
  attach boxed title to top left={yshift=-2mm,xshift=6mm},
  boxed title style={colback=alertbdr!25,colframe=alertbdr,arc=3pt,boxrule=0.6pt}}

\newtcolorbox{warnbox}[1][Obrigação Contratual]{enhanced,breakable,
  colback=warnbg,colframe=warnbdr,boxrule=0.8pt,arc=4pt,
  left=10pt,right=10pt,top=5pt,bottom=5pt,
  title=\textbf{#1},fonttitle=\small\bfseries\color{warnbdr},
  attach boxed title to top left={yshift=-2mm,xshift=6mm},
  boxed title style={colback=warnbg,colframe=warnbdr,arc=3pt,boxrule=0.6pt}}

\newtcolorbox{infobox}[1][Informação]{enhanced,breakable,
  colback=infobg,colframe=infobdr,boxrule=0.8pt,arc=4pt,
  left=10pt,right=10pt,top=5pt,bottom=5pt,
  title=\textbf{#1},fonttitle=\small\bfseries\color{infobdr},
  attach boxed title to top left={yshift=-2mm,xshift=6mm},
  boxed title style={colback=infobg,colframe=infobdr,arc=3pt,boxrule=0.6pt}}

\newtcolorbox{conclusabox}{enhanced,
  colback=navy!8,colframe=navy,boxrule=1pt,arc=5pt,
  left=12pt,right=12pt,top=8pt,bottom=8pt}

% ── Títulos ──────────────────────────────────────────────────────────
\usepackage{titlesec}
\titleformat{\section}{\color{navy}\large\bfseries}{\thesection.}{0.6em}{}[\color{navy}\titlerule]
\titlespacing{\section}{0pt}{22pt}{8pt}
\titleformat{\subsection}{\color{medblue}\normalsize\bfseries}{\thesubsection}{0.6em}{}
\titlespacing{\subsection}{0pt}{14pt}{5pt}
\titleformat{\subsubsection}{\color{navy!80}\normalsize\bfseries\itshape}{\thesubsubsection}{0.5em}{}
\titlespacing{\subsubsection}{0pt}{10pt}{4pt}

% ── Tabelas ──────────────────────────────────────────────────────────
\usepackage{booktabs,array,tabularx,longtable,multirow,colortbl,pdflscape}
\renewcommand{\arraystretch}{1.40}
\newcolumntype{C}[1]{>{\centering\arraybackslash}p{#1}}
\newcolumntype{R}[1]{>{\raggedleft\arraybackslash}p{#1}}
\newcolumntype{L}[1]{>{\raggedright\arraybackslash}p{#1}}

% ── Utilitários ──────────────────────────────────────────────────────
\usepackage{float,graphicx,enumitem,setspace,amsmath,caption,chngcntr}
\setlist[itemize]{leftmargin=1.8em,itemsep=2pt,topsep=3pt,parsep=0pt}
\setlist[enumerate]{leftmargin=1.8em,itemsep=2pt,topsep=3pt,parsep=0pt}
\captionsetup{font=small,labelfont=bf,skip=4pt}
\counterwithout{footnote}{section}
\usepackage[colorlinks=true,linkcolor=navy,urlcolor=medblue,pdfborder={0 0 0}]{hyperref}

% ── Macros ───────────────────────────────────────────────────────────
\newcommand{\art}[1]{art.\,#1}
\newcommand{\pct}{\%}
\newcommand{\Rs}{R\$\,}
\newcommand{\assinatura}[3]{%
  \begin{center}\vspace{0.2cm}
  \rule{7cm}{0.5pt}\\[3pt]
  \textbf{#1}\\\small #2\\\small Mat.: \underline{\hspace{2cm}}\\\small #3
  \end{center}}

% ── Dados do processo ─────────────────────────────────────────────────
\newcommand{\trNum}{TR-SMS/BORBA-001/2025}
\newcommand{\etpRef}{ETP-SMS/BORBA-001/2025}
\newcommand{\dfdRef}{DFD-SMS/BORBA-001/2025}
\newcommand{\procNum}{[NÚMERO DO PROCESSO SEI/E-DOC]}
\newcommand{\exercicio}{2025}

% ═════════════════════════════════════════════════════════════════════
\begin{document}
% ═════════════════════════════════════════════════════════════════════

% ── Capa ─────────────────────────────────────────────────────────────
\begin{titlepage}
  \pagecolor{navy}\color{white}\setlength{\parindent}{0pt}\centering
  \vspace*{1cm}
  {\small\bfseries\color{ltblue}PREFEITURA MUNICIPAL DE BORBA}\\[2pt]
  {\small\color{ltblue!70}SECRETARIA MUNICIPAL DE SAÚDE}\\[0.4cm]
  \textcolor{ltblue!50}{\rule{0.85\linewidth}{0.8pt}}\\[0.6cm]
  {\fontsize{11}{14}\selectfont\bfseries\color{ltblue!80}
    TERMO DE REFERÊNCIA --- TR}\\[0.4cm]
  {\fontsize{22}{28}\selectfont\bfseries
    Aquisição de Materiais de\\[4pt]Limpeza e Higiene}\\[0.25cm]
  {\fontsize{13}{16}\selectfont\color{ltblue!80}
    Pregão Eletrônico --- Ata de Registro de Preços}\\[0.6cm]
  \textcolor{ltblue!50}{\rule{0.85\linewidth}{0.8pt}}\\[1cm]
  \begin{tabular}{@{}r@{\hspace{8pt}}p{0.56\linewidth}@{}}
    \color{ltblue!70}\small\bfseries N.\textsuperscript{o} do TR: &
      \small \trNum \\[4pt]
    \color{ltblue!70}\small\bfseries DFD de origem: &
      \small \dfdRef \\[4pt]
    \color{ltblue!70}\small\bfseries ETP de referência: &
      \small \etpRef \\[4pt]
    \color{ltblue!70}\small\bfseries Processo: &
      \small \procNum \\[4pt]
    \color{ltblue!70}\small\bfseries Modalidade: &
      \small Pregão Eletrônico (\art{17}, I) --- ARP (\art{82}) \\[4pt]
    \color{ltblue!70}\small\bfseries Critério: &
      \small Menor Preço por lote (\art{33}, I) \\[4pt]
    \color{ltblue!70}\small\bfseries Valor Estimado: &
      \small \Rs 2.000.000,00 (dois milhões de reais) \\[4pt]
    \color{ltblue!70}\small\bfseries Vigência da ARP: &
      \small 12 meses, prorrogável (\art{84}, \S\,1.\textsuperscript{o}) \\[4pt]
    \color{ltblue!70}\small\bfseries Base Legal: &
      \small Lei n.\textsuperscript{o} 14.133/2021; Decreto n.\textsuperscript{o} 11.462/2023 \\[4pt]
  \end{tabular}
  \vfill
  \textcolor{ltblue!40}{\rule{0.85\linewidth}{0.4pt}}\\[0.3cm]
  {\tiny\color{ltblue!50}Documento elaborado em conformidade com o
    \art{6.\textsuperscript{o}}, XXIII, e \art{40} da Lei n.\textsuperscript{o}
    14.133/2021 --- Borba/AM, \exercicio}
\end{titlepage}
\nopagecolor
\setlength{\parindent}{1.2cm}

% ── Ficha de identificação ────────────────────────────────────────────
\begin{table}[H]\centering\small
\caption*{\textbf{\color{navy}FICHA DE IDENTIFICAÇÃO DO TERMO DE REFERÊNCIA}}
\begin{tabular}{>{\bfseries\color{navy}}p{4.8cm} L{9.2cm}}
\rowcolor{navy!10}\multicolumn{2}{l}{\textbf{\color{navy}Dados do Documento}} \\
\midrule
N.\textsuperscript{o} do TR:     & \trNum \\
\rowcolor{rowalt}
DFD de origem:                    & \dfdRef \\
ETP de referência:                & \etpRef \\
\rowcolor{rowalt}
N.\textsuperscript{o} do Processo: & \procNum \\
Unidade Demandante:               & Secretaria Municipal de Saúde de Borba/AM \\
\rowcolor{rowalt}
Objeto:                           & Aquisição de materiais de limpeza e higiene
                                    para a Rede de Atenção à Saúde de Borba/AM,
                                    mediante Ata de Registro de Preços \\
Modalidade:                       & Pregão Eletrônico \\
\rowcolor{rowalt}
Critério de julgamento:           & Menor Preço por lote (\art{33}, I) \\
Valor Estimado Global:            & \Rs 2.000.000,00 \\
\rowcolor{rowalt}
Vigência da ARP:                  & 12 meses, prorrogável (\art{84}, \S\,1.\textsuperscript{o}) \\
Número de lotes:                  & 6 \\
\rowcolor{rowalt}
Responsável pela elaboração:      & [NOME, CARGO, MATRÍCULA] \\
Data de elaboração:               & Borba/AM, \today \\
\bottomrule
\end{tabular}
\end{table}

\vspace{0.4cm}
\tableofcontents
\clearpage

% ═════════════════════════════════════════════════════════════════════
\section{Objeto}
% ═════════════════════════════════════════════════════════════════════

\begin{legalbox}[\art{6.\textsuperscript{o}}, XXIII, ``a'', Lei n.\textsuperscript{o} 14.133/2021]
\small O TR deve conter a \textbf{descrição do objeto da contratação}, bem como a
fundamentação da necessidade, os requisitos da contratação e os critérios de
mensuração dos resultados.
\end{legalbox}

\subsection{Descrição do Objeto}

Aquisição de \textbf{materiais de limpeza e higiene} --- compreendendo saneantes de
uso hospitalar e doméstico, antissépticos, equipamentos de proteção individual (EPI)
de limpeza, sacos para acondicionamento de resíduos de serviços de saúde (RSS) e
utensílios de higienização --- para abastecimento contínuo da \textbf{Rede de Atenção
à Saúde do Município de Borba/AM}, composta por 9 Unidades Básicas de Saúde e o
Hospital Municipal Vó Mundoca (40 leitos), por meio de \textbf{Ata de Registro de
Preços (ARP)}, nos termos do \art{82} da Lei n.\textsuperscript{o}~14.133/2021.

\subsection{Fundamentação e Referência ao ETP}

A necessidade da contratação, a justificativa da solução adotada, o levantamento de
mercado, a estimativa de quantidades e o mapeamento de riscos encontram-se
documentados no \textbf{Estudo Técnico Preliminar \etpRef}, integrante dos autos do
processo \procNum.

% ═════════════════════════════════════════════════════════════════════
\section{Descrição dos Lotes e Itens}
% ═════════════════════════════════════════════════════════════════════

\begin{legalbox}[\art{40}, \S\,1.\textsuperscript{o}, I e II, Lei n.\textsuperscript{o} 14.133/2021]
\small O objeto é dividido em \textbf{lotes por afinidade técnica}, sendo vedado o
agrupamento artificial que restrinja a competição. Cada lote constitui unidade
autônoma para fins de habilitação e julgamento de propostas.
\end{legalbox}

\subsection{Estrutura de Lotes}

\begin{table}[H]\centering
\caption{Estrutura de lotes e valores estimados}
\small\rowcolors{2}{rowalt}{white}
\begin{tabular}{C{1.2cm} L{6.0cm} C{1.5cm} C{2.8cm} C{1.8cm}}
  \rowcolor{navy}
  \color{white}\textbf{Lote} &
  \color{white}\textbf{Descrição} &
  \color{white}\textbf{Classe} &
  \color{white}\textbf{Valor Est. (R\$)} &
  \color{white}\textbf{N.\textsuperscript{o} Itens} \\
  \toprule
  01 & Álcoois e antissépticos topológicos                       & A & 820.000,00 &  4 \\
  02 & Saneantes hospitalares (enzimático, glutaraldeído, hipoclorito) & A & 480.000,00 & 3 \\
  03 & EPI de limpeza e materiais descartáveis                   & A & 300.000,00 &  3 \\
  04 & Sacos para resíduos de saúde e comuns                     & B & 180.000,00 &  2 \\
  05 & Desinfetantes, sabonetes e complementares                 & B & 120.000,00 &  3 \\
  06 & Utensílios de limpeza e consumíveis gerais                & C & 100.000,00 &  6 \\
  \rowcolor{ltblue}
  \multicolumn{3}{c}{\textbf{TOTAL GERAL}} &
  \textbf{2.000.000,00} & \textbf{21} \\
  \bottomrule
\end{tabular}
\end{table}

\subsection{Especificação Técnica Detalhada dos Itens}

\begin{landscape}
\begin{longtable}{C{0.5cm} L{4.2cm} C{1.0cm} C{1.5cm} C{1.8cm} C{2.2cm} L{5.5cm}}
\caption{Especificação técnica completa dos 21 itens} \label{tab:itens} \\

\rowcolor{navy}
\color{white}\small\textbf{It.} &
\color{white}\small\textbf{Descrição} &
\color{white}\small\textbf{Un.} &
\color{white}\small\textbf{Qtde.\newline Mensal} &
\color{white}\small\textbf{Qtde.\newline Anual} &
\color{white}\small\textbf{Pr. Ref.\newline Unit. (R\$)} &
\color{white}\small\textbf{Especificação Técnica Mínima} \\
\toprule
\endfirsthead
\multicolumn{7}{l}{\footnotesize\textit{(continuação)}} \\
\rowcolor{navy}
\color{white}\small\textbf{It.} &
\color{white}\small\textbf{Descrição} &
\color{white}\small\textbf{Un.} &
\color{white}\small\textbf{Qtde.\newline Mensal} &
\color{white}\small\textbf{Qtde.\newline Anual} &
\color{white}\small\textbf{Pr. Ref.\newline Unit. (R\$)} &
\color{white}\small\textbf{Especificação Técnica Mínima} \\
\toprule
\endhead
\bottomrule\endfoot
\bottomrule\endlastfoot

\rowcolor{navy!8}\multicolumn{7}{l}{\textbf{\small Lote 01 --- Álcoois e Antissépticos}} \\
\rowcolor{rowalt}
01 & Álcool Etílico Hidratado 70\% --- frasco 500\,mL &
Fr & 1.680 & 20.160 & 7,50 &
Teor alcoólico 69,3--70,5\,\pct\,v/v; com ou sem glicerina; embalagem PEAD; registro ANVISA; prazo validade $\geq$ 18\,meses; tampa dosadora. \\

02 & Álcool Etílico Hidratado 70\% --- galão 5\,L &
Gl & 168 & 2.016 & 36,00 &
Mesmas especificações do Item 01; embalagem galão PEAD; lacre de segurança. \\

\rowcolor{rowalt}
03 & Álcool Isopropílico 70\% --- frasco 1\,L &
Fr & 25 & 300 & 19,00 &
Grau técnico; pureza $\geq$ 99,5\,\pct; embalagem PEAD; registro ANVISA; uso para superfícies e equipamentos eletrônicos; tampa com vedante. \\

04 & Sabonete Líquido Antisséptico Clorexidina 2\% --- 1\,L &
Fr & 90 & 1.080 & 18,50 &
Digluconato de clorexidina 2\,\pct; pH 5,0--7,0; registro ANVISA como produto de higiene pessoal grau 2; embalagem com bomba dosadora; prazo validade $\geq$ 24\,meses. \\

\rowcolor{navy!8}\multicolumn{7}{l}{\textbf{\small Lote 02 --- Saneantes Hospitalares}} \\
\rowcolor{rowalt}
05 & Detergente Enzimático Concentrado --- frasco 1\,L &
Fr & 30 & 360 & 42,94 &
Formulação multienzimática (protease, amilase, lipase); diluição mínima 1:200; pH neutro (6,5--7,5); registro ANVISA; biodegradável; não corrosivo em instrumentos cirúrgicos; prazo validade $\geq$ 24\,meses. \\

06 & Hipoclorito de Sódio 1\,\pct\ --- embalagem 5\,L &
Fr & 280 & 3.360 & 5,80 &
Cloro ativo 0,9--1,1\,\pct; pH 11,0--13,0; registro ANVISA; embalagem opaca PEAD; tampa rosca com lacre; prazo validade $\geq$ 12\,meses; FISPQ disponível. \\

\rowcolor{rowalt}
07 & Glutaraldeído 2\,\pct\ Tamponado --- frasco 1\,L &
Fr & 12 & 144 & 65,00 &
Concentração 2\,\pct $\pm$ 0,1\,\pct; pH 7,5--8,5; uso como desinfetante de nível alto e esterilizante; registro ANVISA; prazo $\geq$ 24\,meses; instrução de preparo de solução ativada. \\

\rowcolor{navy!8}\multicolumn{7}{l}{\textbf{\small Lote 03 --- EPI de Limpeza e Descartáveis}} \\
08 & Luva de Procedimento Látex Não Estéril --- caixa 100\,un.\ --- tam. M &
Cx & 80 & 960 & 38,00 &
Látex natural; não estéril; ambidestra; tam. M; textura na ponta dos dedos; espessura $\geq$ 0,10\,mm; registro ANVISA; ABNT NBR 13.392; AQL $\leq$ 1,5. \\

\rowcolor{rowalt}
09 & Avental Descartável TNT --- caixa 50\,un. &
Cx & 30 & 360 & 55,00 &
TNT gramatura $\geq$ 30\,g/m²; manga longa; fechamento nas costas; tamanho único; cor branca; registro ANVISA; embalagem individual não necessária. \\

10 & Papel Toalha Branco Interfolhado 2 dobras --- pacote 1.000\,folhas &
Pct & 36 & 432 & 22,00 &
Celulose virgem $\geq$ 100\,\pct; 2 dobras; folha 21$\times$22\,cm; alvura $\geq$ 82\,ISO; carga $\leq$ 0,5\,g/m²; sem perfume; pacote plástico. \\

\rowcolor{navy!8}\multicolumn{7}{l}{\textbf{\small Lote 04 --- Sacos para Resíduos}} \\
\rowcolor{rowalt}
11 & Saco p/ Lixo Infectante Vermelho 100\,L --- rolo 100\,un. &
Rl & 60 & 720 & 95,00 &
Cor vermelha; PEAD reciclado; espessura $\geq$ 0,10\,mm; capacidade 100\,L; símbolo de resíduo infectante impresso; ABNT NBR 9.191; limite $\leq$ 1\,kg de carga. \\

12 & Saco p/ Lixo Doméstico Preto 100\,L --- rolo 100\,un. &
Rl & 40 & 480 & 42,00 &
Cor preta; PEBD/PEAD; espessura $\geq$ 0,08\,mm; capacidade 100\,L; resistência $\geq$ 8\,kgf (norma ABNT NBR 9.191); rolo contínuo. \\

\rowcolor{navy!8}\multicolumn{7}{l}{\textbf{\small Lote 05 --- Desinfetantes e Complementares}} \\
13 & Desinfetante Quaternário de Amônio --- embalagem 5\,L &
Fr & 45 & 540 & 28,00 &
Cloreto de benzalcônio $\geq$ 8\,\pct\ ou equivalente; pH 7,0--8,5; registro ANVISA; para superfícies fixas hospitalares; fragrância neutra; biodegradável; diluição $\geq$ 1:80. \\

\rowcolor{rowalt}
14 & Sabão em Pó Multiuso --- embalagem 1\,kg &
Pct & 50 & 600 & 9,50 &
Tensoativo aniônico; alvejante óptico; enzimas multiação; poder de limpeza $\geq$ 90\,\pct\ (ASTM D-4008); embalagem plástica com tampa hermética. \\

15 & Detergente Líquido Neutro Concentrado --- frasco 500\,mL &
Fr & 60 & 720 & 3,80 &
pH 6,5--7,5; concentrado; rendimento $\geq$ 80 lavagens/frasco; registro ANVISA (se aplicável); biodegradável $\geq$ 90\,\pct; tampa com vedante de segurança. \\

\rowcolor{navy!8}\multicolumn{7}{l}{\textbf{\small Lote 06 --- Utensílios e Consumíveis Gerais}} \\
\rowcolor{rowalt}
16 & Água Sanitária --- embalagem 2\,L &
Fr & 120 & 1.440 & 5,20 &
Hipoclorito de sódio 2,0--2,5\,\pct\ de cloro ativo; registro ANVISA; embalagem PEAD opaca; prazo validade $\geq$ 12\,meses; FISPQ disponível. \\

17 & Vassoura Nylon Profissional &
Un & 12 & 144 & 18,00 &
Cerdas nylon 6/66; comprimento cerdas $\geq$ 7\,cm; base polipropileno; encaixe rosca $\frac{3}{4}''$; resistente a saneantes; cabo não incluso. \\

\rowcolor{rowalt}
18 & Rodo 60\,cm --- Cabo Alumínio &
Un & 10 & 120 & 22,00 &
Borracha dupla face; lâmina 60\,cm; cabo alumínio $\emptyset$ 25\,mm; rosca universal; comprimento mínimo 1,40\,m; resistente a saneantes; borracha $\geq$ 4\,mm espessura. \\

19 & Pano de Chão Algodão --- 55$\times$75\,cm &
Un & 30 & 360 & 7,50 &
100\,\pct\ algodão alvejado; trama $\geq$ 32 fios/cm; resistência tração $\geq$ 200\,N; lavável; sem corante que manche superfícies. \\

\rowcolor{rowalt}
20 & Balde Plástico 20\,L c/ Alça Metálica &
Un & 5 & 60 & 25,00 &
Polipropileno alimentício; capacidade 20\,L; resistência $\geq$ 40\,kg carga; alça metálica com revestimento; resistente a saneantes; escala em relevo. \\

21 & Esponja Dupla Face Verde/Amarela &
Un & 40 & 480 & 3,50 &
Face abrasiva não metálica (fibra verde); face macia (espuma de poliuretano); dimensões $\geq$ 110$\times$75$\times$25\,mm; resistente a detergentes; embalagem individual. \\

\end{longtable}
\end{landscape}

\begin{alertbox}[Marca e Equivalência]
\small
As especificações acima \textbf{não indicam marca}. Serão aceitos produtos equivalentes
que atendam os parâmetros técnicos mínimos. A comprovação de equivalência é de
responsabilidade do licitante, mediante apresentação de laudo de laboratório acreditado
pelo INMETRO ou documentação técnica do fabricante, a ser analisada pelo fiscal técnico
do contrato/ARP.
\end{alertbox}

% ═════════════════════════════════════════════════════════════════════
\section{Requisitos da Contratação}
% ═════════════════════════════════════════════════════════════════════

\begin{legalbox}[\art{6.\textsuperscript{o}}, XXIII, ``d''; \art{40}, I, Lei n.\textsuperscript{o} 14.133/2021]
\small O TR deve conter os \textbf{requisitos da contratação}, incluindo requisitos
técnicos, legais e de sustentabilidade que o contratado deverá cumprir.
\end{legalbox}

\subsection{Requisitos Técnicos Obrigatórios}

\begin{enumerate}
  \item \textbf{Registro ANVISA}: todos os saneantes dos Lotes 01 a 05 devem possuir
    registro, notificação ou isenção de registro na ANVISA, vigentes na data da entrega,
    nos termos da Lei n.\textsuperscript{o}~6.360/1976 e RDC ANVISA
    n.\textsuperscript{o}~59/2021;

  \item \textbf{Ficha de Dados de Segurança (FISPQ)}: o fornecedor deverá entregar,
    juntamente com a primeira remessa de cada produto químico, a FISPQ atualizada,
    conforme ABNT NBR 14.725;

  \item \textbf{Prazo de validade mínimo}: os produtos devem ser entregues com prazo
    de validade residual de \textbf{no mínimo 12 meses} contados da data de entrega,
    salvo especificação diferente no Item~2.2;

  \item \textbf{Integridade das embalagens}: embalagens violadas, amassadas, com vazamento
    ou sem rótulo legível serão recusadas no ato do recebimento, sem ônus ao Município;

  \item \textbf{Rastreabilidade}: os produtos devem apresentar número de lote e data
    de fabricação impressos na embalagem, permitindo rastreabilidade em caso de recall
    pela ANVISA.
\end{enumerate}

\subsection{Requisitos Logísticos}

\begin{enumerate}
  \item \textbf{Preço CIF Borba/AM}: as propostas devem ser apresentadas em preço
    \textbf{CIF Borba} (\textit{Cost, Insurance and Freight} --- Custo, Seguro e Frete),
    incluindo todos os custos de transporte desde a origem até o almoxarifado da SMS,
    vedada a oferta em preço FOB Manaus ou ex-works;

  \item \textbf{Logística fluvial}: o fornecedor deve comprovar, na fase de habilitação,
    capacidade logística para entrega em Borba/AM por via fluvial (trecho Manaus--Borba,
    Rio Madeira), mediante carta de empresa de navegação ou contrato de fretamento;

  \item \textbf{Prazos de entrega}: conforme Seção~5 deste TR;

  \item \textbf{Acondicionamento}: os produtos devem ser entregues em caixas de papelão
    reforçado ou engradados plásticos, protegidos de umidade e impacto, adequados ao
    transporte fluvial de longa distância.
\end{enumerate}

\subsection{Requisitos de Sustentabilidade}

Nos termos do \art{11}, IV, da Lei n.\textsuperscript{o}~14.133/2021:

\begin{enumerate}
  \item \textbf{Biodegradabilidade}: para os saneantes dos Lotes 01 a 05, será dada
    \textbf{preferência} a produtos com laudo de biodegradabilidade emitido por
    laboratório acreditado pelo INMETRO (RDC ANVISA n.\textsuperscript{o}~59/2021,
    \art{6.\textsuperscript{o}});
  \item \textbf{Embalagens}: preferência por embalagens com menor volume de resíduos,
    passíveis de reciclagem ou refis;
  \item \textbf{Logística reversa}: o fornecedor se obriga a recolher embalagens vazias
    de produtos perigosos (solventes, saneantes ácidos/alcalinos), nos termos do
    \art{33} da Lei n.\textsuperscript{o}~12.305/2010 (PNRS).
\end{enumerate}

% ═════════════════════════════════════════════════════════════════════
\section{Modelo de Execução do Objeto}
% ═════════════════════════════════════════════════════════════════════

\begin{legalbox}[\art{6.\textsuperscript{o}}, XXIII, ``e''; \art{40}, \S\,1.\textsuperscript{o}, III, Lei n.\textsuperscript{o} 14.133/2021]
\small O TR deve conter o \textbf{modelo de execução do objeto}, descrevendo como o
contrato/ARP será executado, os prazos, as condições de entrega e as obrigações das partes.
\end{legalbox}

\subsection{Instrumento de Contratação}

O fornecimento será formalizado por meio de \textbf{Ata de Registro de Preços (ARP)},
nos termos do \art{82} da Lei n.\textsuperscript{o}~14.133/2021 e do Decreto
n.\textsuperscript{o}~11.462/2023, com as seguintes características:

\begin{table}[H]\centering\small
\begin{tabular}{>{\bfseries\color{navy}}p{4.8cm} L{9.2cm}}
\rowcolor{navy!8}\multicolumn{2}{l}{\textbf{\color{navy}Parâmetros da ARP}} \\
\midrule
Vigência:           & 12 (doze) meses a contar da data de publicação, prorrogável
                      por igual período (\art{84}, \S\,1.\textsuperscript{o}) \\
\rowcolor{rowalt}
Natureza:           & Não obriga o Município à aquisição integral (\art{83}) \\
Instrumento:        & Ordem de Fornecimento (OF) emitida pelo órgão gerenciador \\
\rowcolor{rowalt}
Quantidades:        & Teto máximo por item; sem mínimo garantido ao fornecedor \\
Adesão (carona):    & Vedada sem autorização expressa do órgão gerenciador (\art{86}, \S\,4.\textsuperscript{o}) \\
\rowcolor{rowalt}
Reequilíbrio:       & Mediante solicitação fundamentada, com base em índice IPCA/INPC (\art{124}) \\
\bottomrule
\end{tabular}
\end{table}

\subsection{Forma de Fornecimento e Ordens de Fornecimento}

\begin{enumerate}
  \item O fornecimento será \textbf{parcelado}, mediante emissão de \textbf{Ordens de
    Fornecimento (OF)} ao longo da vigência da ARP;

  \item Cada OF especificará os itens, as quantidades, o prazo de entrega e o local de
    recebimento;

  \item As OF serão emitidas com \textbf{antecedência mínima de 20 dias} em relação à
    data prevista de entrega, respeitando os prazos de logística fluvial
    Manaus--Borba;

  \item O intervalo mínimo entre OF para o mesmo item é de \textbf{30 dias}, salvo
    situação de emergência documentada pelo fiscal do contrato;

  \item O Município não está obrigado a emitir OF que totalize menos de
    \Rs 5.000,00 por entrega, para garantir viabilidade econômica do frete fluvial;

  \item As quantidades solicitadas em cada OF serão calibradas ao \textbf{ciclo
    hidrológico do Rio Madeira}:
    \begin{itemize}
      \item Período de cheia (jan--jun): OF com cobertura de 30--45 dias;
      \item Período de estiagem (jul--out): OF com cobertura de 45--60 dias
        (Estoque de Reserva Estratégica);
      \item Período de transição (nov--dez): OF com cobertura de 30 dias.
    \end{itemize}
\end{enumerate}

\subsection{Local de Entrega}

Todos os itens deverão ser entregues no \textbf{Almoxarifado Central da Secretaria
Municipal de Saúde de Borba/AM}, sito à [ENDEREÇO COMPLETO --- RUA, N.\textsuperscript{o},
BAIRRO, CEP], Borba/AM, em dias úteis, no horário das 08:00 às 12:00 horas e das
14:00 às 17:00 horas, mediante agendamento prévio com o fiscal do contrato.

\subsection{Prazos de Entrega}

\begin{table}[H]\centering
\caption{Prazos máximos de entrega por lote}
\small\rowcolors{2}{rowalt}{white}
\begin{tabular}{C{1.2cm} L{6.0cm} C{3.0cm} C{3.0cm}}
  \rowcolor{navy}
  \color{white}\textbf{Lote} &
  \color{white}\textbf{Descrição} &
  \color{white}\textbf{Prazo Cheia (dias)} &
  \color{white}\textbf{Prazo Seca (dias)} \\
  \toprule
  01--02 & Álcoois, antissépticos e saneantes hospitalares & 25 & 35 \\
  03     & EPI de limpeza e descartáveis                    & 25 & 35 \\
  04--05 & Resíduos, desinfetantes e complementares         & 30 & 40 \\
  06     & Utensílios e consumíveis gerais                  & 30 & 40 \\
  \bottomrule
\end{tabular}
\end{table}

\noindent Os prazos são contados em \textbf{dias corridos} a partir do \textbf{recebimento
da Ordem de Fornecimento}. O período de seca do Rio Madeira compreende os meses de
julho a outubro de cada ano.

\begin{warnbox}[Consequência do Atraso na Entrega]
\small
O descumprimento do prazo de entrega sujeita o fornecedor à multa de mora prevista
na Seção~10 deste TR, sem prejuízo da possibilidade de rescisão unilateral da ARP por
inexecução parcial, nos termos do \art{137} da Lei n.\textsuperscript{o}~14.133/2021.
\end{warnbox}

% ═════════════════════════════════════════════════════════════════════
\section{Recebimento e Aceitação}
% ═════════════════════════════════════════════════════════════════════

\begin{legalbox}[\art{140}, Lei n.\textsuperscript{o} 14.133/2021]
\small O recebimento do objeto observará o disposto no \art{140}, com recebimento
provisório pelo servidor designado e recebimento definitivo por servidor ou comissão
designada pela autoridade competente.
\end{legalbox}

\subsection{Recebimento Provisório}

O \textbf{servidor responsável pelo almoxarifado} realizará o recebimento provisório
no ato da entrega, verificando:

\begin{enumerate}
  \item Conferência quantitativa (quantidade entregue vs.\ OF emitida);
  \item Integridade das embalagens (sem violação, amassamento ou vazamento);
  \item Identificação dos produtos (rótulo com nome, lote, fabricação e validade);
  \item Prazo de validade residual mínimo exigido (Seção~3.1, item 3);
  \item Temperatura de transporte, quando aplicável (álcoois: $<$ 30\,\textdegree C).
\end{enumerate}

O recebimento provisório será registrado em \textbf{Termo de Recebimento Provisório
(TRP)}, emitido em até \textbf{2 dias úteis} após a entrega.

\subsection{Recebimento Definitivo}

O \textbf{fiscal técnico do contrato} realizará o recebimento definitivo em até
\textbf{10 dias úteis} após o provisório, mediante:

\begin{enumerate}
  \item Verificação da conformidade técnica com as especificações do TR;
  \item Confirmação do registro ANVISA vigente (consulta ao portal ANVISA em tempo real);
  \item Análise da FISPQ (Lotes 01 a 05, primeira entrega e a cada alteração de lote);
  \item Aprovação de laudo de equivalência, se aplicável (Seção~2.3, \textit{alertbox});
  \item Emissão do \textbf{Termo de Recebimento Definitivo (TRD)} e liberação para pagamento.
\end{enumerate}

\begin{alertbox}[Recusa do Objeto]
\small
Produtos que não atendam às especificações serão \textbf{devolvidos ao fornecedor}
mediante Notificação de Não Conformidade (NNC), com prazo de \textbf{5 dias úteis}
para substituição. A devolução não desobriga o fornecedor do prazo original de
entrega, sujeitando-o às penalidades cabíveis.
\end{alertbox}

% ═════════════════════════════════════════════════════════════════════
\section{Modelo de Gestão do Contrato e da ARP}
% ═════════════════════════════════════════════════════════════════════

\begin{legalbox}[\art{6.\textsuperscript{o}}, XXIII, ``f''; arts. 117--120, Lei n.\textsuperscript{o} 14.133/2021]
\small O TR deve conter o \textbf{modelo de gestão do contrato/ARP}, incluindo a
designação do gestor e do fiscal e os mecanismos de controle e avaliação da execução.
\end{legalbox}

\subsection{Gestor da ARP}

A SMS designará formalmente um \textbf{Gestor da ARP}, servidor efetivo, com publicação
em Diário Oficial, com as seguintes atribuições principais (\art{117},
\S\,1.\textsuperscript{o}):

\begin{itemize}
  \item Coordenar a emissão das Ordens de Fornecimento;
  \item Controlar o saldo físico-financeiro da ARP;
  \item Decidir sobre pedidos de reequilíbrio econômico-financeiro;
  \item Propor a prorrogação, rescisão ou revogação da ARP;
  \item Comunicar ao fornecedor registrado qualquer alteração nas condições da ARP.
\end{itemize}

\subsection{Fiscal Técnico da ARP}

A SMS designará um \textbf{Fiscal Técnico}, servidor com conhecimento técnico em
limpeza hospitalar e/ou logística de insumos, com as seguintes atribuições (\art{117},
\S\,1.\textsuperscript{o}):

\begin{itemize}
  \item Verificar a conformidade técnica de cada entrega;
  \item Emitir Termos de Recebimento Provisório e Definitivo;
  \item Registrar inconformidades e aplicar Notificações de Não Conformidade;
  \item Monitorar os KPIs de desempenho (Seção~8);
  \item Instruir processos de penalidade, quando necessário.
\end{itemize}

\subsection{Indicadores de Gestão (KPIs)}

\begin{table}[H]\centering
\caption{Indicadores de desempenho da ARP}
\small\rowcolors{2}{rowalt}{white}
\begin{tabular}{L{4.0cm} L{3.8cm} C{2.2cm} C{2.2cm}}
  \rowcolor{navy}
  \color{white}\textbf{Indicador} &
  \color{white}\textbf{Fórmula} &
  \color{white}\textbf{Meta} &
  \color{white}\textbf{Frequência} \\
  \toprule
  Taxa de Ruptura de Estoque   & Dias em falta / Dias do período   & $<\,0{,}1\,\pct$      & Mensal \\
  Custo por m² Higienizado     & Custo total / 5.150\,m²            & $\leq$ \Rs 8,50/m²   & Trimestral \\
  Aderência ao Ponto de Pedido & OF no prazo / Total de OF          & $>\,95\,\pct$         & Mensal \\
  Perda por Validade           & Valor descartado / Valor adquirido & $<\,2\,\pct$          & Semestral \\
  Cobertura do ES              & $E_{\text{atual}}/ES_{\text{calc}}$ & $>\,100\,\pct$       & Quinzenal \\
  Prazo médio de entrega       & $\sum$ (prazo real) / N OF         & $\leq$ prazo contrat. & Mensal \\
  \bottomrule
\end{tabular}
\end{table}

% ═════════════════════════════════════════════════════════════════════
\section{Critérios de Medição e Pagamento}
% ═════════════════════════════════════════════════════════════════════

\begin{legalbox}[\art{6.\textsuperscript{o}}, XXIII, ``g''; \art{141}, Lei n.\textsuperscript{o} 14.133/2021]
\small O TR deve conter os \textbf{critérios de medição e pagamento} da parcela
correspondente ao fornecimento efetivamente realizado.
\end{legalbox}

\subsection{Faturamento e Nota Fiscal}

\begin{enumerate}
  \item O pagamento será realizado por \textbf{Ordem de Fornecimento} --- cada OF
    gerará uma nota fiscal correspondente;
  \item A \textbf{Nota Fiscal Eletrônica (NF-e)} deverá conter: número da OF, número
    da ARP, CNPJ do fornecedor, discriminação dos itens com quantidades, preços
    unitários e valores totais, e o campo de observações com o número do processo
    administrativo;
  \item A NF-e deverá ser emitida \textbf{após a emissão do Termo de Recebimento
    Definitivo (TRD)}, vedado o faturamento de itens ainda não recebidos definitivamente.
\end{enumerate}

\subsection{Prazo e Forma de Pagamento}

\begin{enumerate}
  \item O pagamento será realizado em \textbf{até 30 dias corridos} contados da
    apresentação da NF-e ao setor de finanças da SMS, após a emissão do TRD;
  \item O pagamento será efetuado por \textbf{transferência bancária} (TED/PIX),
    para conta corrente em nome do fornecedor, previamente cadastrada;
  \item São condições \textbf{suspensivas} do pagamento:
    \begin{itemize}
      \item Certidão Negativa de Débitos Previdenciários (CND/INSS) válida;
      \item Certificado de Regularidade do FGTS (CRF) válido;
      \item Certidão de Débitos Relativos a Tributos Federais válida;
      \item Ausência de penalidade ativa que impeça a contratação.
    \end{itemize}
  \item Ocorrendo \textbf{atraso no pagamento por culpa do Município}, incidirão
    sobre o valor devido correção monetária pelo IPCA e juros de mora de 0,5\,\pct\ ao mês,
    conforme \art{141}, \S\,3.\textsuperscript{o}.
\end{enumerate}

\subsection{Glosas e Descontos}

Serão descontados do pagamento:

\begin{itemize}
  \item Valores correspondentes a itens devolvidos por não conformidade;
  \item Multas de mora já aplicadas e não pagas voluntariamente;
  \item Diferenças de preço em caso de entrega de item em quantidade superior à OF,
    quando não autorizada previamente pelo gestor.
\end{itemize}

% ═════════════════════════════════════════════════════════════════════
\section{Habilitação}
% ═════════════════════════════════════════════════════════════════════

\begin{legalbox}[arts. 62 a 70, Lei n.\textsuperscript{o} 14.133/2021]
\small Os requisitos de habilitação devem ser os estritamente \textbf{indispensáveis}
para garantir o cumprimento das obrigações contratuais, vedada a imposição de
exigências desnecessárias (\art{62}, \S\,1.\textsuperscript{o}).
\end{legalbox}

\subsection{Habilitação Jurídica (\art{66})}

\begin{itemize}
  \item Registro comercial, no caso de empresa individual;
  \item Ato constitutivo, estatuto ou contrato social em vigor, devidamente registrado
    na Junta Comercial ou cartório competente;
  \item Decreto de autorização, em se tratando de empresa ou sociedade estrangeira.
\end{itemize}

\subsection{Habilitação Fiscal, Social e Trabalhista (arts. 68 e 69)}

\begin{itemize}
  \item CNPJ ativo na Receita Federal;
  \item Certidão Negativa de Débitos Federais (RFB/PGFN);
  \item Certidão Negativa Estadual e Municipal do domicílio do licitante;
  \item Certificado de Regularidade do FGTS (CRF --- Caixa Econômica Federal);
  \item Certidão Negativa de Débitos Trabalhistas (CNDT --- TST).
\end{itemize}

\subsection{Habilitação Técnica (\art{67})}

\begin{enumerate}
  \item \textbf{Autorização de Funcionamento ANVISA} para distribuição de saneantes
    (Lotes 01 a 05), nos termos da Lei n.\textsuperscript{o}~6.360/1976 e Portaria
    SVS n.\textsuperscript{o}~2.814/1998;

  \item \textbf{Comprovação de logística fluvial}: carta de empresa de navegação
    habilitada para o trecho Manaus--Borba/AM, ou contrato de fretamento com empresa
    de transporte fluvial autorizada pela ANTAQ, com prazo compatível com a vigência
    da ARP;

  \item \textbf{Atestado de Capacidade Técnica}: atestado(s) fornecido(s) por pessoa
    jurídica de direito público ou privado, comprovando fornecimento anterior de
    material de limpeza e higiene em quantidade equivalente a no mínimo \textbf{50\,\pct\
    do quantitativo anual} do respectivo lote (\art{67}, I);

  \item \textbf{Declaração de Amostras}: o licitante primeiro colocado deverá
    apresentar amostras dos itens de cada lote em que sagrarem-se vencedores, no
    prazo de \textbf{5 dias úteis} após o encerramento da fase de lances, para
    avaliação técnica pelo fiscal.
\end{enumerate}

\subsection{Habilitação Econômico-Financeira (\art{69})}

\begin{itemize}
  \item Capital social ou patrimônio líquido mínimo de \textbf{10\,\pct\ do valor
    estimado do respectivo lote} (\art{69}, I);
  \item Certidão negativa de falência ou recuperação judicial, expedida pelo
    distribuidor da sede do licitante;
  \item Balanço patrimonial referente ao último exercício social (dispensável para
    empresas constituídas no mesmo exercício da licitação).
\end{itemize}

\subsection{Benefícios para ME/EPP}

Aplicam-se os benefícios dos arts. 44 e 45 da Lei Complementar n.\textsuperscript{o}
123/2006:

\begin{itemize}
  \item \textbf{Empate ficto}: propostas de ME/EPP até 5\,\pct\ superiores à melhor
    proposta serão consideradas empatadas, com direito de preferência de desempate;
  \item \textbf{Regularização fiscal}: prazo de 5 dias úteis para regularização de
    eventuais restrições fiscais, após declarada a vencedora.
\end{itemize}

% ═════════════════════════════════════════════════════════════════════
\section{Critérios de Seleção do Fornecedor e Julgamento}
% ═════════════════════════════════════════════════════════════════════

\begin{legalbox}[\art{6.\textsuperscript{o}}, XXIII, ``h''; \art{33}, I, Lei n.\textsuperscript{o} 14.133/2021]
\small O critério de julgamento é o \textbf{menor preço} por lote, nos termos do
\art{33}, I, sendo este o mais adequado para bens com especificação técnica precisa.
\end{legalbox}

\subsection{Critério de Julgamento}

\begin{itemize}
  \item Critério: \textbf{Menor Preço Global por Lote};
  \item Cada lote é independente para fins de julgamento e habilitação;
  \item É admitida a participação de um mesmo licitante em mais de um lote;
  \item Não é admitida a adjudicação parcial dentro do mesmo lote (o vencedor
    de cada lote fornece todos os itens do lote);
  \item O preço de referência por item constante da Tabela~\ref{tab:itens} serve
    como limite máximo de aceitabilidade para cada item: propostas com preço unitário
    superior ao valor de referência em mais de \textbf{10\,\pct} serão desclassificadas.
\end{itemize}

\subsection{Proposta de Preços}

A proposta deverá conter:

\begin{enumerate}
  \item Preço unitário e total para cada item do lote, em reais, com duas casas decimais;
  \item Preços \textbf{CIF Borba/AM}, já incluídos todos os custos de frete, seguro,
    impostos e encargos;
  \item Prazo de validade da proposta: mínimo 60 dias;
  \item Marca, modelo e fabricante para cada item ofertado;
  \item Número(s) de registro ANVISA dos produtos dos Lotes 01 a 05.
\end{enumerate}

% ═════════════════════════════════════════════════════════════════════
\section{Obrigações do Fornecedor Registrado}
% ═════════════════════════════════════════════════════════════════════

\begin{enumerate}
  \item Manter, durante toda a vigência da ARP, as mesmas condições de habilitação
    exigidas na licitação;
  \item Cumprir os prazos de entrega estabelecidos nas Ordens de Fornecimento;
  \item Comunicar ao gestor da ARP, com antecedência mínima de \textbf{15 dias},
    qualquer dificuldade logística que possa comprometer o cumprimento de uma OF;
  \item Substituir, às suas expensas, no prazo de \textbf{5 dias úteis}, produtos
    devolvidos por não conformidade;
  \item Manter vigentes todos os registros ANVISA dos itens registrados na ARP; em
    caso de cancelamento, comunicar imediatamente o gestor e providenciar produto
    substituto equivalente, sujeito a nova análise técnica;
  \item Apresentar FISPQ atualizada sempre que houver alteração na formulação dos
    produtos;
  \item Responsabilizar-se pelo transporte, acondicionamento e seguro dos produtos
    até o almoxarifado da SMS;
  \item Permitir e facilitar a realização de \textbf{auditoria} nas instalações de
    armazenagem e distribuição pelo fiscal ou por órgão de controle competente;
  \item Não subcontratar o fornecimento sem prévia autorização escrita do gestor;
  \item Emitir nota fiscal eletrônica conforme os requisitos da Seção~7.1.
\end{enumerate}

% ═════════════════════════════════════════════════════════════════════
\section{Obrigações do Município (Órgão Gerenciador)}
% ═════════════════════════════════════════════════════════════════════

\begin{enumerate}
  \item Emitir as Ordens de Fornecimento com antecedência mínima de 20 dias;
  \item Efetuar o pagamento no prazo estipulado, após emissão do TRD;
  \item Designar formalmente gestor e fiscal da ARP, com publicação em Diário Oficial,
    no ato da assinatura;
  \item Fornecer ao fornecedor todas as informações necessárias ao cumprimento
    das obrigações contratuais;
  \item Notificar formalmente o fornecedor em caso de não conformidade ou atraso,
    antes de aplicar qualquer sanção;
  \item Manter o almoxarifado da SMS em condições adequadas para o recebimento
    e armazenagem dos produtos, conforme ABNT NBR 14.725 e NR-32.
\end{enumerate}

% ═════════════════════════════════════════════════════════════════════
\section{Sanções Administrativas}
% ═════════════════════════════════════════════════════════════════════

\begin{legalbox}[arts. 155 a 163, Lei n.\textsuperscript{o} 14.133/2021]
\small As sanções serão aplicadas conforme o \art{156}, observado o contraditório
e a ampla defesa (\art{157}).
\end{legalbox}

\begin{table}[H]\centering
\caption{Tabela de sanções administrativas}
\small\rowcolors{2}{rowalt}{white}
\begin{tabular}{L{4.5cm} C{2.0cm} L{7.2cm}}
  \rowcolor{navy}
  \color{white}\textbf{Infração} &
  \color{white}\textbf{Sanção} &
  \color{white}\textbf{Parâmetro} \\
  \toprule
  Atraso na entrega (por dia) &
    Multa moratória &
    0,5\,\pct\ ao dia sobre o valor da OF, até o 30.\textsuperscript{o} dia \\
  Atraso na entrega $>$ 30 dias &
    Multa + rescisão &
    Multa compensatória de 10\,\pct\ do valor total da ARP (\art{162}) \\
  Entrega de produto não conforme &
    Advertência + multa &
    Advertência na 1.\textsuperscript{a} ocorrência; multa de 5\,\pct\ da OF nas
    ocorrências seguintes \\
  Não substituição do produto devolvido no prazo &
    Multa &
    2\,\pct\ ao dia sobre o valor do item, até o 5.\textsuperscript{o} dia útil \\
  Entrega de produto sem registro ANVISA &
    Rescisão imediata + multa compensatória &
    10\,\pct\ do valor total da ARP; comunicação à ANVISA \\
  Reincidência em qualquer infração &
    Impedimento de licitar e contratar &
    Até 3 anos, nos termos do \art{156}, III \\
  \bottomrule
\end{tabular}
\end{table}

\begin{warnbox}[Defesa Prévia]
\small
Antes de aplicada qualquer penalidade, será garantido ao fornecedor \textbf{prazo de
10 dias úteis para apresentação de defesa} (\art{157}). O não exercício do direito
de defesa no prazo implicará preclusão.
\end{warnbox}

% ═════════════════════════════════════════════════════════════════════
\section{Reajuste e Revisão de Preços}
% ═════════════════════════════════════════════════════════════════════

\begin{legalbox}[\art{92}, I; \art{124}, I, Lei n.\textsuperscript{o} 14.133/2021]
\small Os preços registrados na ARP são passíveis de revisão (\textit{reequilíbrio
econômico-financeiro}) e de reajuste periódico, conforme condições estabelecidas.
\end{legalbox}

\subsection{Reajuste Anual}

\begin{itemize}
  \item Os preços serão reajustados anualmente com base no \textbf{IPCA acumulado}
    nos últimos 12 meses, a contar da data do orçamento de referência da pesquisa
    de preços;
  \item O reajuste será processado de ofício pelo gestor, sem necessidade de
    requerimento do fornecedor;
  \item Somente incidirá para ARP prorrogada além de 12 meses.
\end{itemize}

\subsection{Reequilíbrio Econômico-Financeiro}

\begin{itemize}
  \item O fornecedor poderá solicitar revisão de preços mediante fato imprevisível e
    extraordinário que altere o equilíbrio da equação econômico-financeira, nos termos
    do \art{124}, I (\textit{áleas econômicas extraordinárias});
  \item A solicitação deverá ser fundamentada com documentação comprobatória
    (notas fiscais de insumos, índices setoriais, etc.);
  \item O gestor terá \textbf{30 dias} para analisar e decidir o pedido;
  \item Em caso de recusa, o fornecedor poderá solicitar o cancelamento de seu
    registro, sem penalidade, nos termos do \art{84}, \S\,2.\textsuperscript{o}.
\end{itemize}

% ═════════════════════════════════════════════════════════════════════
\section{Estimativa de Preços de Referência}
% ═════════════════════════════════════════════════════════════════════

\begin{legalbox}[\art{23}; \art{40}, \S\,2.\textsuperscript{o}, I, Lei n.\textsuperscript{o} 14.133/2021]
\small A estimativa de preços de referência foi obtida por pesquisa de mercado, conforme
a IN SEGES/MGI n.\textsuperscript{o}~65/2021.
\end{legalbox}

\subsection{Metodologia}

\begin{enumerate}
  \item \textbf{Painel de Preços} (SEGES/MGI) --- compras homologadas nos últimos
    12 meses para itens idênticos ou similares;
  \item \textbf{BPS --- Banco de Preços em Saúde} (Ministério da Saúde) --- para
    itens de uso hospitalar;
  \item \textbf{3 cotações diretas} junto a distribuidores atacadistas de Manaus;
  \item \textbf{Contratos anteriores} da SMS Borba (2022--2024).
\end{enumerate}

\noindent O valor de referência de cada item corresponde à \textbf{média aritmética}
das fontes acima, \textbf{majorado em 13\,\pct}\ a título de frete fluvial
Manaus--Borba e perdas de transporte (custo CIF), conforme memória de cálculo
juntada ao processo.

\subsection{Resumo do Valor Estimado}

\begin{table}[H]\centering
\caption{Valor estimado global por lote (CIF Borba/AM)}
\small\rowcolors{2}{rowalt}{white}
\begin{tabular}{C{1.2cm} L{6.5cm} R{3.0cm} R{2.5cm}}
  \rowcolor{navy}
  \color{white}\textbf{Lote} &
  \color{white}\textbf{Descrição} &
  \color{white}\textbf{Valor Est. 12 meses (R\$)} &
  \color{white}\textbf{\% do Total} \\
  \toprule
  01 & Álcoois e antissépticos topológicos                  & 820.000,00 & 41,0\,\pct \\
  02 & Saneantes hospitalares                               & 480.000,00 & 24,0\,\pct \\
  03 & EPI de limpeza e descartáveis                        & 300.000,00 & 15,0\,\pct \\
  04 & Sacos para resíduos                                  & 180.000,00 &  9,0\,\pct \\
  05 & Desinfetantes, sabonetes e complementares            & 120.000,00 &  6,0\,\pct \\
  06 & Utensílios de limpeza e consumíveis                  & 100.000,00 &  5,0\,\pct \\
  \rowcolor{ltblue}
  \multicolumn{2}{c}{\textbf{TOTAL GERAL (CIF Borba/AM)}} & \textbf{2.000.000,00} & \textbf{100\,\pct} \\
  \bottomrule
\end{tabular}
\end{table}

\begin{infobox}[Sigilosidade do Orçamento Referencial]
\small
O detalhamento dos preços unitários de referência por item \textbf{não será publicado
no edital}, nos termos do \art{24} da Lei n.\textsuperscript{o}~14.133/2021. O valor
global estimado por lote será divulgado. O orçamento detalhado integra o processo
como documento de acesso restrito até o encerramento da fase de lances.
\end{infobox}

% ═════════════════════════════════════════════════════════════════════
\section{Adequação Orçamentária}
% ═════════════════════════════════════════════════════════════════════

\begin{legalbox}[\art{40}, \S\,1.\textsuperscript{o}, IV; \art{11}, IV, Lei n.\textsuperscript{o} 14.133/2021]
\small
\end{legalbox}

\begin{table}[H]\centering\small
\begin{tabular}{>{\bfseries\color{navy}}p{4.8cm} L{9.2cm}}
\rowcolor{navy!8}\multicolumn{2}{l}{\textbf{\color{navy}Adequação Orçamentária}} \\
\midrule
Unidade Orçamentária: & [CÓDIGO UO] --- Fundo Municipal de Saúde \\
\rowcolor{rowalt}
Programa de Trabalho: & [CÓDIGO PROGRAMA] \\
Ação:                 & [CÓDIGO AÇÃO] --- Manutenção da Rede de Atenção à Saúde \\
\rowcolor{rowalt}
Natureza da Despesa:  & 33.90.30 --- Material de Consumo \\
Fonte de Recurso:     & [FONTE --- Rec. Próprio / FMS / SUS] \\
\rowcolor{rowalt}
Dotação Disponível:   & \Rs 2.000.000,00 (confirmada em [DATA]) \\
\bottomrule
\end{tabular}
\end{table}

\noindent A despesa relativa ao fornecimento parcelado via ARP será empenhada por
ocasião de cada Ordem de Fornecimento, na forma do \art{58} da Lei
n.\textsuperscript{o}~4.320/1964 e \art{82}, \S\,4.\textsuperscript{o} da Lei
n.\textsuperscript{o}~14.133/2021.

% ═════════════════════════════════════════════════════════════════════
\section{Disposições Finais}
% ═════════════════════════════════════════════════════════════════════

\begin{enumerate}
  \item Os casos omissos serão resolvidos com base na Lei n.\textsuperscript{o}
    14.133/2021 e, subsidiariamente, nos princípios gerais do direito;
  \item O \textbf{foro competente} para dirimir questões oriundas deste TR e da ARP
    é o da Comarca de \textbf{Borba/AM}, com renúncia expressa a qualquer outro;
  \item A ARP poderá ser \textbf{cancelada} pelo gestor nas hipóteses do
    \art{84}, \S\,2.\textsuperscript{o}: preços registrados tornem-se superiores ao
    mercado; ocorrer fato superveniente que justifique; fornecedor recusar redução dos
    preços nos termos do \art{84}, \S\,3.\textsuperscript{o};
  \item Este TR poderá ser \textbf{atualizado} pelo gestor nos casos em que a
    evolução do mercado ou a experiência acumulada na execução indicarem necessidade
    de ajustes, para fins de prorrogação da ARP;
  \item O presente instrumento integra o processo administrativo \procNum\ e é
    parte indissociável do edital de licitação, do qual reproduzirá os elementos
    essenciais.
\end{enumerate}

% ── Aprovação ─────────────────────────────────────────────────────────
\vspace{1.5cm}

\begin{conclusabox}
\begin{center}
  \textbf{\color{navy}\large APROVAÇÃO DO TERMO DE REFERÊNCIA}
\end{center}
\smallskip
Declaro que o presente Termo de Referência foi elaborado com base em critérios
técnicos objetivos, em conformidade com o \art{6.\textsuperscript{o}}, XXIII, da
Lei n.\textsuperscript{o} 14.133/2021, sem direcionamento para fornecedor
específico, e que as especificações adotadas são as estritamente necessárias para
atender às necessidades da Secretaria Municipal de Saúde de Borba/AM.
\end{conclusabox}

\vspace{1.4cm}
\begin{minipage}{0.47\linewidth}
\assinatura{[Nome e Cargo]}{Responsável pela elaboração do TR}{SMS Borba/AM}
\end{minipage}
\hfill
\begin{minipage}{0.47\linewidth}
\assinatura{[Nome e Cargo]}{Secretário(a) Municipal de Saúde}{Aprovação hierárquica}
\end{minipage}

\vspace{1cm}
\begin{minipage}{0.47\linewidth}
\assinatura{[Nome e Cargo]}{Procurador(a)-Geral do Município}{Aprovação jurídica (PGM)}
\end{minipage}
\hfill
\begin{minipage}{0.47\linewidth}
\assinatura{[Nome e Cargo]}{Prefeito(a) Municipal de Borba}{Autoridade competente}
\end{minipage}

\vspace{0.8cm}\noindent\hfill{\small Borba/AM, \today}

\vspace{1.2cm}
\noindent\textcolor{navy!30}{\rule{\linewidth}{0.5pt}}

\noindent{\footnotesize\itshape
\textbf{Referências Normativas:}
Lei n.\textsuperscript{o} 14.133/2021 (NLLCA) \textbullet{}
Decreto n.\textsuperscript{o} 11.462/2023 (SRP) \textbullet{}
IN SEGES/MGI n.\textsuperscript{o} 65/2021 (Pesquisa de Preços) \textbullet{}
IN SEGES/MGI n.\textsuperscript{o} 58/2022 (PCA e DFD) \textbullet{}
RDC ANVISA n.\textsuperscript{o} 59/2021 (Saneantes) \textbullet{}
RDC ANVISA n.\textsuperscript{o} 15/2012 (Artigos Hospitalares) \textbullet{}
RDC ANVISA n.\textsuperscript{o} 216/2004 \textbullet{}
Lei n.\textsuperscript{o} 6.360/1976 (Vigilância Sanitária) \textbullet{}
NR-32 (MTE) \textbullet{}
LC n.\textsuperscript{o} 123/2006 (ME e EPP) \textbullet{}
Lei n.\textsuperscript{o} 12.305/2010 (PNRS) \textbullet{}
Lei n.\textsuperscript{o} 4.320/1964 \textbullet{}
ABNT NBR 14.725 \textbullet{} ABNT NBR 9.191 \textbullet{} ABNT NBR 13.392.}

\end{document}
